\documentclass[12pt,a4paper]{article}
\usepackage{fontspec}
\setmainfont{Times New Roman}
\usepackage[top=2.0cm,bottom=2.0cm,left=2.5cm,right=2.0cm]{geometry}
\usepackage{enumitem}
\usepackage{tabularx}
\usepackage{xltabular}
\usepackage{booktabs}
\usepackage{array}
\newcolumntype{Y}{>{\centering\arraybackslash}X}
\newcolumntype{C}[1]{>{\centering\arraybackslash}p{#1}}
\newcolumntype{L}[1]{>{\raggedright\arraybackslash}p{#1}}
\newcolumntype{R}[1]{>{\raggedleft\arraybackslash}p{#1}}
\pagestyle{plain}
\linespread{1.15}
\setlength{\parskip}{4pt plus 1pt minus 1pt}
\sloppy

\begin{document}
% --- HƯỚNG DẪN ĐIỀU TRA ---
\begin{center}
    \fontsize{12pt}{14pt}\selectfont
    \textbf{TỔNG CỤC CẢNH SÁT — BAN ĐIỀU PHỐI HỒ SƠ VỤ ÁN}\\[-0.2em]
    \rule{6cm}{0.6pt}\\[0.5cm]
    \textbf{\Large CẨM NANG HƯỚNG DẪN ĐIỀU TRA \& NGUYÊN TẮC PHÁ ÁN}\\[0.15cm]
    \textit{}
\end{center}
\vspace{0.3cm}



\vspace{0.3cm}\noindent\textbf{\large I. HƯỚNG DẪN ĐỐI VỚI BẰNG CHỨNG VẬT LÝ}\\[0.15cm]


\begin{itemize}[leftmargin=*, itemsep=2pt, topsep=2pt, parsep=0pt]
  \item \textbf{Mở hồ sơ theo thứ tự quy định:}
\begin{itemize}[leftmargin=*, itemsep=2pt, topsep=2pt, parsep=0pt]
  \item \textbf{Bộ hồ sơ ban đầu:} Luôn được mở ra đầu tiên khi bắt đầu vụ  án để nắm bắt báo cáo tiếp nhận nguồn tin, hiện trường ban đầu và pháp y sơ bộ.
  \item \textbf{Các bộ hồ sơ kế tiếp (Ký hiệu A, B, C):} Tuyệt đối chỉ được mở theo từng giai đoạn điều tra khi có chỉ dẫn.

\end{itemize}
  \item \textbf{Quy tắc ghi nhận \& Điền số thứ tự (STT):}
\begin{itemize}[leftmargin=*, itemsep=2pt, topsep=2pt, parsep=0pt]
  \item Mỗi tài liệu, vật chứng đều được gán số thứ tự (STT) cụ thể \textit{(ví dụ: P1, P2, P3... hoặc EV-...)}.
  \item Khi nhập đáp án chứng minh suy luận trên hệ thống, thám tử cần điền chính xác STT của bằng chứng tương ứng để hệ thống xác thực dữ liệu.


\end{itemize}
\end{itemize}

\vspace{0.3cm}\noindent\textbf{\large II. HƯỚNG DẪN SỬ DỤNG HỆ THỐNG WEBSITE}\\[0.15cm]



\vspace{0.2cm}\noindent\textbf{1. Bảng điều khiển chính}\\[0.1cm]


\begin{itemize}[leftmargin=*, itemsep=2pt, topsep=2pt, parsep=0pt]
  \item \textbf{Nghi phạm:} Chọn đối tượng cụ thể để tập trung xác minh. Đối tượng chỉ đủ điều kiện đưa vào diện nghi phạm khi thỏa mãn đồng thời hai yếu tố:
\begin{itemize}[leftmargin=*, itemsep=2pt, topsep=2pt, parsep=0pt]
  \item \textbf{Có động cơ:} Tồn tại mâu thuẫn về tài sản, tình cảm, trả thù hoặc có lý do trực tiếp để ra tay với nạn nhân.
  \item \textbf{Ngoại phạm bất hợp lý:} Xuất hiện sự sai lệch, mâu thuẫn giữa lời khai của đối tượng với dữ liệu khách quan thực tế.

\end{itemize}
  \item \textbf{Mở rộng điều tra:}
\begin{itemize}[leftmargin=*, itemsep=2pt, topsep=2pt, parsep=0pt]
  \item \textbf{Bổ sung chứng cứ:} Trả lời các câu hỏi logic nghiệp vụ để thu thập thêm manh mối, mở khóa tài liệu mới.
  \item \textbf{Khám xét lại hiện trường:} Nhằm thu thập thêm các manh mối bị bỏ sót khi những chứng cứ hiện tại chưa đủ kết luận hoặc vụ án phát sinh tình tiết mới. Thao tác này chỉ được thực hiện khi đạt đủ điều kiện kích hoạt lệnh \textit{(ấn chọn trên website để xem chi tiết yêu cầu)}.
  \item \textbf{Đề nghị truy tố:} Thực hiện khi đã tự tin vào kết luận cuối cùng. Kết luận truy tố bắt buộc phải tích chọn chính xác các chứng cứ chứng minh, không chấp nhận suy đoán thiếu căn cứ.

\end{itemize}
\end{itemize}

\vspace{0.2cm}\noindent\textbf{2. Bảng điều khiển phụ}\\[0.1cm]


\begin{itemize}[leftmargin=*, itemsep=2pt, topsep=2pt, parsep=0pt]
  \item \textbf{Điện thoại nạn nhân:} Truy cập trực tiếp vào thiết bị của nạn nhân để kiểm tra tin nhắn, cuộc gọi, ghi chú và dữ liệu số.
  \item \textbf{Gợi ý phá án:} Cung cấp hướng đi khi bế tắc, theo đúng tiến trình điều tra tại thời điểm đó. Hãy kiên nhẫn suy luận trước khi sử dụng.
  \item \textbf{Chơi lại:} Khôi phục toàn bộ tiến trình điều tra về trạng thái ban đầu.
\end{itemize}
\end{document}
